\documentclass[11pt]{article}
\usepackage[margin=1in]{geometry}
\usepackage[T1]{fontenc}
\usepackage{textcomp}
\usepackage{hyperref}
\usepackage{listings}
\usepackage{enumitem}

\title{Chroma Native SMD Design}
\author{}
\date{\today}

\begin{document}
\maketitle

\section{Overview}
This document specifies a native Stochastic Molecular Dynamics (SMD) implementation
in Chroma. The design follows the paper \href{https://arxiv.org/pdf/1105.4749}{arXiv:1105.4749} and \href{https://luscher.web.cern.ch/luscher/notes/smd-ergodicity.pdf}{Luscher's notes}. The implementation is orthogonal to existing HMC
code: new modules and a new executable are added, while changes to existing files
are minimal and limited to build and registration glue.

\section{Goals}
\begin{itemize}[leftmargin=*]
  \item Provide an SMD driver that uses existing Chroma monomials, actions, and integrators.
  \item Support three variants:
    \begin{itemize}
      \item GHMC-style SMD with accept/reject.
      \item SMD without accept/reject but with action evaluation (for diagnostics).
      \item SMD without accept/reject and without action evaluation (fast mode).
    \end{itemize}
  \item Provide a new XML-driven executable (\texttt{smd}) similar to \texttt{hmc}.
  \item Keep integrators unchanged (Omelyan, minnorm2/4, force-gradient).
\end{itemize}


\section{Algorithm Summary}
The SMD trajectory is composed of a partial momentum refresh, molecular dynamics
integration, and optional accept/reject. The algorithm follows the common pattern:
\begin{enumerate}
  \item Refresh momenta with friction parameter \(\gamma\) and effective step size \(\Delta t\).
  \item Refresh pseudo-fermions (as in HMC initialization).
  \item Integrate the equations of motion using a selected symplectic integrator.
  \item Optionally accept/reject based on \(\Delta H\).
\end{enumerate}

The refresh step uses a standard Ornstein--Uhlenbeck update with an effective
step size \(\Delta t\) set to the outermost trajectory length:
\begin{equation}
  P \leftarrow c_1 P + c_2 \eta, \quad c_1 = e^{-\gamma\Delta t},\quad c_2 = \sqrt{1-c_1^2}
\end{equation}
where \(\eta\) is a Gaussian momentum field.
In Chroma, \(\Delta t\) is taken as the outermost integrator \(\tau\) (the total
MD time for the update). For nested integrators, sub-integrator step sizes are
internal to the MD evolution and do not enter the refresh.

\section{Mapping to Chroma Components}
\subsection{Existing Components Reused}
\begin{itemize}[leftmargin=*]
  \item Monomials and actions: existing HMC monomial factory and action registry.
  \item Integrators: existing LCM integrators (Omelyan, minnorm2/4, force-gradient), as well as nested integrators for multi-time scale integration
  \item RNG and lattice field types: QDP++ and Chroma utilities.
  \item Configuration I/O: existing gauge config reader/writer and XML output helpers.
\end{itemize}


\subsection{New Components}
\begin{itemize}[leftmargin=*]
  \item \texttt{update/molecdyn/smd}: SMD driver and params.
  \item \texttt{mainprogs/main/smd.cc}: new executable, XML-driven.
  \item \texttt{py\_chroma} binding: run SMD from Python without XML.
\end{itemize}

\section{Module and File Layout}
\begin{itemize}[leftmargin=*]
  \item \texttt{src/chroma/lib/update/molecdyn/smd/smd.h}
  \item \texttt{src/chroma/lib/update/molecdyn/smd/smd.cc}
  \item \texttt{src/chroma/lib/update/molecdyn/smd/smd\_params.h}
  \item \texttt{src/chroma/lib/update/molecdyn/smd/smd\_params.cc}
  \item \texttt{src/chroma/mainprogs/main/smd.cc}
  \item \texttt{src/chroma/mainprogs/main/compare\_gauge.cc}
\end{itemize}

Only build and registration glue should touch existing files.

\section{SMD Parameters}
 A minimal parameter set includes:
\begin{itemize}[leftmargin=*]
  \item \texttt{gamma}: friction parameter (>0).
  \item \texttt{AcceptReject}: accept/reject on (\texttt{True}) or off (\texttt{False}).
  \item \texttt{actions}: list of action indices (as in HMC).
  \item \texttt{monomials}: standard Chroma monomial list.
  \item \texttt{integrator}: standard Chroma integrator parameters with explicit
        step counts (\texttt{n\_steps}) and trajectory length (\texttt{tau}).
\end{itemize}

\section{Execution Flow}
\subsection{High-Level Flow}
\begin{enumerate}
  \item Load gauge configuration, momenta and set gauge BC/state.
  	\item if momenta file is not present initialize with gaussian.
  \item Build monomials and integrator from XML.
  \item Initialize SMD state (momenta and pseudo-fermions).
  \item For each update:
    \begin{enumerate}
      \item Refresh momenta (and pseudo-fermions if configured).
      \item Integrate MD equations for one trajectory segment.
      \item Optional accept/reject; if rejected, restore previous fields and flip momenta.
      \item Save configuration and  momenta if needed and log XML.
    \end{enumerate}
\end{enumerate}

\subsection{Accept/Reject Variants}
\begin{itemize}[leftmargin=*]
  \item \textbf{Full SMD (GHMC):} accept/reject on each update (\texttt{AcceptReject=True}).
  \item \textbf{No-acc with actions:} skip accept/reject, still compute action start/end.
  \item \textbf{No-acc fast:} skip accept/reject and skip action evaluation.
\end{itemize}

\section{XML Interface (Executable)}
The \texttt{smd} executable should mirror \texttt{hmc} structure and reuse existing
blocks for configuration I/O and measurements. A skeletal XML layout:
\begin{lstlisting}[basicstyle=\ttfamily\small]
<smd>
  <UpdateParams>
    <n_updates>10</n_updates>
    <save_interval>1</save_interval>
    <start_update_num>0</start_update_num>
    <cfg_file>...</cfg_file>
    <cfg_type>...</cfg_type>
  </UpdateParams>

  <SMDParams>
    <gamma>0.3</gamma>
    <AcceptReject>True</AcceptReject>
  </SMDParams>

  <Monomials> ... </Monomials>
  <MDIntegrator> ... </MDIntegrator>

  <Measurements> ... optional ... </Measurements>
</smd>
\end{lstlisting}


\section{Data Ownership and State}
\begin{itemize}[leftmargin=*]
  \item The SMD driver owns transient momentum and pseudo-fermion fields.
  \item Gauge field is managed by existing Chroma gauge state.
  \item On reject: restore gauge field and momenta; flip momentum sign.
  \item On accept: renormalize gauge links and leave field updated.
\end{itemize}

\section{Logging and Diagnostics}
\begin{itemize}[leftmargin=*]
  \item Report start/end action contributions and \(\Delta H\).
  \item Report acceptance rate when \texttt{AcceptReject=True}.
  \item Log integrator statistics as in HMC (forces, solver stats).
\end{itemize}



\end{document}
